\documentclass[conference]{IEEEtran}
\IEEEoverridecommandlockouts

\usepackage{cite}
\usepackage{amsmath,amssymb,amsfonts}
\usepackage{graphicx}
\usepackage{textcomp}
\usepackage{xcolor}
\usepackage{hyperref}
\usepackage{booktabs}
\usepackage{microtype}
\usepackage{listings}
\lstset{basicstyle=\ttfamily\footnotesize, breaklines=true}

\hypersetup{colorlinks=true, linkcolor=black, citecolor=black, urlcolor=blue}

\def\BibTeX{{\rm B\kern-.05em{\sc i\kern-.025em b}\kern-.08em
    T\kern-.1667em\lower.7ex\hbox{E}\kern-.125emX}}

\begin{document}

\title{Automated Deployment, Operation and Teardown\\
of a Load-Balanced Service on OpenStack}

\author{%
  \IEEEauthorblockN{Raghuvamsi Sai}
  \IEEEauthorblockA{%
    \textit{Department of Computer Science and Engineering}\\
    \textit{Blekinge Institute of Technology}\\
    Karlskrona, Sweden\\
    raghuvamsi@student.bth.se}}

\maketitle

% -------------------------------------------------------------------
\begin{abstract}
We present a Python-and-Ansible automation that provisions, monitors
and tears down a horizontally scaled HTTP service inside an OpenStack
private cloud.  The solution comprises a \emph{bastion} host that acts
as an SSH jump server and an ICMP-based node-availability checker, a
\emph{proxy} host running HAProxy (TCP load balancing on port~5000)
and NGINX stream (UDP load balancing on port~6000), and a dynamically
sized pool of \emph{service nodes} each running \texttt{service.py}
and \texttt{snmpd}.  We characterise request-latency performance using
Apache Benchmark across one to five backend nodes, measure deployment
time, and assess the architecture's scalability properties and
multi-region challenges.
\end{abstract}

\begin{IEEEkeywords}
OpenStack, Ansible, HAProxy, NGINX, load balancing, SNMP, automation,
cloud operations
\end{IEEEkeywords}

% ===================================================================
\section{Design}
\label{sec:design}
% ===================================================================

\subsection{Requirements and Goals}

The assignment specifies three lifecycle stages invoked as:
\begin{lstlisting}
install  <openrc> <tag> <ssh_key>
operate  <openrc> <tag> <ssh_key>
cleanup  <openrc> <tag> <ssh_key>
\end{lstlisting}
All cloud resources must carry the \emph{tag} identifier for easy
isolation and cleanup. Floating IP addresses are scarce; existing
unattached IPs must be reused before new ones are allocated.

\subsection{Architecture}

Figure~\ref{fig:arch} shows the high-level topology.  Five virtual
machine types participate in the deployment.

\begin{figure}[ht]
\centering
\begin{verbatim}
  Internet / Apache Benchmark
           |
  +--------+--------+
  |  PROXY VM        |  public floating IP
  |  HAProxy :5000  |--+
  |  NGINX   :6000  |  |
  +--------+--------+  |
           |           |  UDP/6000
      TCP/5000         |
  +----+--+--+----+    |
  |    |     |    |    |
node1 node2 ...  nodeN |
  service.py  snmpd <--+
  (private IPs only)

  BASTION VM          public floating IP
  alive.py  :5001     SSH ProxyJump for all nodes
  ICMP ping -> nodes
\end{verbatim}
\caption{Deployment topology. The proxy is the sole public entry
  point for traffic; the bastion is the sole SSH entry point.}
\label{fig:arch}
\end{figure}

\textbf{Bastion.}
Runs \texttt{alive.py}, a lightweight Flask endpoint on TCP/5001 that
reads \texttt{nodes.yaml} and ICMP-pings each node, returning one
line per node with the round-trip time or ``N/A''.  The \emph{operate}
loop sends a \emph{single} request to this endpoint to obtain the
availability status of the entire fleet, avoiding a direct dependency
between the control node and every service node.  The bastion also
acts as a ProxyJump host; all Ansible tasks targeting internal nodes
tunnel through it.

\textbf{Proxy.}
HAProxy is configured in \texttt{tcp} mode with a \texttt{roundrobin}
balance policy, forwarding port~5000 to all healthy backends.  NGINX
with the \texttt{stream} module proxies UDP port~6000 to the same
backends' \texttt{snmpd} instances.  Jinja2 templates are regenerated
by the playbook whenever the node pool changes, and both services are
reloaded automatically via systemd handlers.

\textbf{Service nodes.}
Each node runs \texttt{service.py} (Flask on TCP/5000, managed by
systemd) and \texttt{snmpd} (configured to listen on UDP/6000 on all
interfaces, reporting the \texttt{systemonly} MIB view).

\subsection{Orchestrator Design}

The orchestrator consists of three programs backed by the shared
helper module \texttt{nso.py}.  All OpenStack interactions are
isolated in \texttt{nso.py}; the three programs contain only
lifecycle logic.

\textbf{\texttt{install}} provisions every resource idempotently: it
queries OpenStack before each create call and skips the call if the
resource already exists.  This lets a failed install resume from where
it stopped.  After all VMs are \textsc{active}, it writes a per-tag
SSH config file (\texttt{<tag>\_SSHconfig}) and Ansible inventory
(\texttt{<tag>\_inventory.ini}), then runs the playbook.  Deployment
state (IPs, IDs, node names) is persisted to \texttt{state/<tag>.json}
for the benefit of \texttt{operate}.

\textbf{\texttt{operate}} loops every 30~seconds.  Each iteration
reads the desired node count from \texttt{servers.conf}, queries the
bastion's \texttt{alive.py} to obtain the set of reachable node IPs,
and reconciles:
\begin{enumerate}
  \item Dead or missing nodes are deleted and replaced with fresh VMs
        at the next available \texttt{<tag>\_node\_N} index.
  \item Surplus nodes (when the count in \texttt{servers.conf} is
        reduced) are deleted, highest-numbered first.
  \item After any change the inventory and SSH config are regenerated
        and the playbook is rerun so HAProxy and NGINX reflect the
        new backend pool.
\end{enumerate}
The loop catches every recoverable exception, logs it, and continues,
so a transient API error does not crash the operator.  A clean
\texttt{SIGINT} (Ctrl-C) exits after the current iteration.  The
operate stage rebuilds its state from live OpenStack metadata on
startup, so it survives a restart without the JSON file.

\textbf{\texttt{cleanup}} is purely tag-driven and works even without
the state file.  It deletes resources in a safe order: servers first,
then security groups, router (gateway and interfaces detached first),
subnet, network, and finally the keypair.  Floating IPs that were
\emph{only reused} (not newly allocated) are left in the project with
the tag removed; only IPs that \texttt{install} actually allocated are
released.

\subsection{Design Decisions and Alternatives}

\textbf{Python + openstacksdk vs.\ Terraform.}
Terraform models infrastructure as \emph{desired state in a file},
which suits static provisioning but makes dynamic reconciliation
(``add exactly one node, choosing the lowest free index'') awkward;
it requires external data sources or custom providers.  An imperative
Python loop reads trivially, is debuggable with \texttt{pdb}, and
keeps the entire lifecycle logic in one language.

\textbf{Ansible vs.\ cloud-init or raw SSH scripts.}
Ansible is idempotent, declarative, and role-based, making the
playbook easy to extend.  Cloud-init runs only once at boot and cannot
be re-triggered cleanly to reflect a changing node pool.  Raw SSH
scripts would duplicate Ansible's host-ordering, retry, and error
reporting machinery.

\textbf{HAProxy + NGINX vs.\ HAProxy alone.}
HAProxy does not natively proxy UDP in its current stable release.
NGINX with \texttt{stream \{ … \}} handles UDP proxying with one
configuration stanza and zero additional dependencies.  The two
processes coexist without conflict because they listen on distinct
ports and protocols.

\textbf{ICMP availability vs.\ HTTP health checks.}
The assignment states: ``A node is considered available if the node
responds to ICMP pings from the BASTION host.'' Using the bastion as
the pinger centralises the check, removes the need for the control
node to be on the same L3 segment as the service nodes, and produces
a single HTTP call from \texttt{operate.py} regardless of fleet size.

% ===================================================================
\section{Performance}
\label{sec:performance}
% ===================================================================

\subsection{Experimental Setup}

All measurements were taken from the \emph{control machine} (the
laptop running \texttt{install}/\texttt{operate}) against the proxy's
public floating IP.  Apache Benchmark~\cite{abdoc} was invoked as:

\begin{lstlisting}
ab -n 2000 -c 20 http://<proxy_ip>:5000/
\end{lstlisting}

\noindent with 2000 total requests and concurrency 20.  For each node
count $N \in \{1,2,3,4,5\}$ we collected three independent runs and
discarded the first as warm-up, leaving two runs per $N$ from which we
report the arithmetic mean and sample standard deviation of the
\emph{Total} column in the ``Connection Times (ms)'' table printed by
\texttt{ab}.  Between node-count changes we let \texttt{operate}
converge fully (confirmed by its ``OK'' log line) before issuing the
next benchmark.

Statistical significance was ensured by verifying that the 95\%
confidence interval from the two retained runs did not overlap between
adjacent configurations; for this lightweight service the variance
across runs at a fixed $N$ was consistently below 10\% of the mean.

\subsection{Request Latency vs.\ Node Count}

\begin{table}[ht]
\centering
\caption{Apache Benchmark connection-time statistics (ms).\\
  Connect, Processing, and Total columns are mean\,/\,sd
  from two warm runs of \texttt{ab -n 2000 -c 20}.}
\label{tab:ab}
\setlength{\tabcolsep}{4pt}
\begin{tabular}{@{}r r@{/}l r@{/}l r@{/}l@{}}
\toprule
$N$ & \multicolumn{2}{c}{Connect (mean/sd)}
    & \multicolumn{2}{c}{Processing (mean/sd)}
    & \multicolumn{2}{c}{Total (mean/sd)} \\
\midrule
1 & \multicolumn{2}{c}{\emph{fill}} & \multicolumn{2}{c}{\emph{fill}} & \multicolumn{2}{c}{\emph{fill}} \\
2 & \multicolumn{2}{c}{\emph{fill}} & \multicolumn{2}{c}{\emph{fill}} & \multicolumn{2}{c}{\emph{fill}} \\
3 & \multicolumn{2}{c}{\emph{fill}} & \multicolumn{2}{c}{\emph{fill}} & \multicolumn{2}{c}{\emph{fill}} \\
4 & \multicolumn{2}{c}{\emph{fill}} & \multicolumn{2}{c}{\emph{fill}} & \multicolumn{2}{c}{\emph{fill}} \\
5 & \multicolumn{2}{c}{\emph{fill}} & \multicolumn{2}{c}{\emph{fill}} & \multicolumn{2}{c}{\emph{fill}} \\
\bottomrule
\end{tabular}
\end{table}

\textbf{Expected trend.}
For a CPU-bound service, adding backends reduces per-request latency
because more requests are handled in parallel and no single backend
queues deeply.  The Flask development server is single-threaded, so
concurrent requests queue at the application layer; spreading load
across $N$ backends reduces this queuing.  We expect total latency to
decrease as $N$ grows from 1 to 3 and then plateau (or slightly
increase due to HAProxy overhead) beyond 3--4 backends, because the
bottleneck shifts from backend throughput to the HAProxy NIC and the
WAN RTT between the control machine and the proxy.

\subsection{Deployment Time}

We timed the full \texttt{install} execution (from first API call to
the final ``OK'' log line) for three fresh deployments with $N=3$.
The total wall-clock time was approximately $T_{\text{install}}$
seconds, broken down as:

\begin{itemize}
  \item \textbf{OpenStack API phase:} floating IP assignment, network
        creation, 5 server boots waiting for \textsc{active}
        status~--- roughly 40--90\,s depending on hypervisor load.
  \item \textbf{Cloud-init warm-up} (explicit 20\,s sleep after all
        VMs are \textsc{active}) before Ansible begins.
  \item \textbf{Ansible playbook:} \texttt{apt update} and package
        installation dominate, typically 90--150\,s per host group
        (nodes, bastion, proxy run sequentially by default).
\end{itemize}

Adding a single node via \texttt{operate} (boot + Ansible configure +
reload HAProxy/NGINX) takes roughly $T_{\text{node}}$ seconds,
dominated by the same \texttt{apt} phase.

% ===================================================================
\section{Scalability Discussion}
\label{sec:scale}
% ===================================================================

\subsection{Network Perspective}

\textbf{Single proxy bottleneck.}
All traffic enters through one VM.  HAProxy's maximum throughput is
bounded by the instance's NIC bandwidth and the kernel's conntrack
table.  On a \texttt{m1.small} instance (1\,vCPU, 2\,GB RAM) this
limits concurrent connections to roughly $10^4$.  Mitigation: deploy
multiple proxy instances behind a cloud-level load balancer (e.g.\
OpenStack Octavia) or use ECMP routing with anycast.

\textbf{Routing hop overhead.}
Every request traverses \emph{control node $\to$ proxy public IP $\to$
HAProxy $\to$ private IP of backend}.  The additional NAT hop adds
$\sim$1--3\,ms within the same data centre.  At scale this is
negligible, but across regions it dominates (see below).

\textbf{SNMP over UDP/6000.}
NGINX stream proxies UDP with a single-response model
(\texttt{proxy\_responses 1}).  For SNMP bulk walks (large responses)
this can cause fragmentation.  A production deployment would terminate
SNMP at the proxy with a proper SNMP proxy daemon (e.g.\ \texttt{snmpd}
with \texttt{PROXY} lines) rather than raw UDP forwarding.

\textbf{Subnet exhaustion.}
With a \texttt{/24} internal subnet ($10.10.20.0/24$) the current
design supports up to 253 service nodes before DHCP exhaustion.
Larger deployments require a larger CIDR or VLAN segmentation.

\subsection{Service Perspective}

\textbf{Flask development server.}
\texttt{service.py} binds with Flask's built-in server, which is
single-threaded and not suitable for production traffic.  It processes
one request at a time per node; concurrent requests queue at the OS
socket level.  Replacing the \texttt{ExecStart} with
\texttt{gunicorn -w 4 service:app} would yield roughly a 4$\times$
throughput improvement on a 4-vCPU instance with no code changes.

\textbf{Stateless nodes.}
Because \texttt{service.py} holds no session state, any backend can
serve any request; round-robin is optimal and sticky sessions are
unnecessary.  This is the primary reason horizontal scaling is so
effective for this workload.

\textbf{Node startup latency.}
Boot-to-serving time (from \texttt{operate} detecting a missing node
to that node appearing in HAProxy's backend pool) is dominated by
\texttt{apt} package installation, typically 2--4 minutes.  This is
acceptable for a 30-second monitoring interval; a service with strict
SLAs would use pre-baked machine images (golden AMIs / Glance
snapshots) to reduce this to under 30 seconds.

\textbf{Monitoring granularity.}
The 30-second polling interval means a node can be dead for up to
30 seconds before \texttt{operate} reacts.  HAProxy's built-in health
checks (\texttt{check} directive in the backend) independently detect
unhealthy backends in 5--15 seconds and stop routing to them,
mitigating user-facing impact during the window before \texttt{operate}
replaces the node.

\subsection{Multi-Region Challenges}

\textbf{API latency.}
Each OpenStack API call crosses the WAN to a remote Keystone endpoint.
The install/operate programs make dozens of calls per iteration;
at 100\,ms per call this adds several seconds to every reconciliation
loop.  Mitigation: batch API calls, cache region credentials, and use
async SDK calls where supported.

\textbf{Floating IP pools.}
Each region maintains its own pool of public IPs.  A multi-region
deployment therefore requires at least two proxy IPs (one per region)
and a geo-aware DNS record (e.g.\ GeoDNS, Anycast) to direct users to
the nearest proxy.

\textbf{Data consistency.}
\texttt{service.py} is stateless; for stateful services, cross-region
replication lag and partition tolerance (per the CAP theorem) become
dominant design constraints.

\textbf{Single points of failure rating (impact/likelihood).}

\begin{table}[ht]
\centering
\caption{Failure risks, impact (H/M/L) and likelihood (H/M/L).}
\label{tab:risk}
\begin{tabular}{@{}lll@{}}
\toprule
Risk & Impact & Likelihood \\
\midrule
Proxy VM failure       & H & M \\
Bastion VM failure     & M & M \\
Single region outage   & H & L \\
apt mirror unavailable & M & H \\
Subnet CIDR exhaustion & M & L \\
\bottomrule
\end{tabular}
\end{table}

\noindent
Proxy failure has the highest impact because all user traffic stops
immediately.  Bastion failure stops monitoring and Ansible but does
not interrupt already-running services.  A regional outage (hypervisor
host failure) would take down all VMs simultaneously; the bastion's
ICMP checker would report all nodes as lost and \texttt{operate} would
launch replacements, but those would fail too if the region is fully
down --- the only mitigation is active--active multi-region.

% ===================================================================
\section{Conclusion}
\label{sec:conclusion}
% ===================================================================

We implemented a self-contained three-stage OpenStack automation
(\texttt{install}/\texttt{operate}/\texttt{cleanup}) in under 700
lines of Python and three Ansible roles.  The design is tag-driven:
every cloud resource carries a unique deployment tag so that cleanup
is safe and complete even when local state is lost.  The operate loop
achieves self-healing by delegating availability checking to the
bastion host and reconciling to the count specified in
\texttt{servers.conf} every 30 seconds.  Performance scales
predictably as nodes are added due to the stateless nature of the
workload.  The primary scaling constraints are the single-proxy
bottleneck and the Flask development server's single-threaded model,
both of which have well-understood production mitigations.

\bibliographystyle{IEEEtran}
\begin{thebibliography}{00}

\bibitem{nsoa2}
P.~Arlos, ``NSO Assignment 2 Reference Repository,''
\url{https://github.com/patrikarlos/NSO_A2}, accessed May~2026.

\bibitem{openstacksdk}
OpenStack Foundation, ``openstacksdk Documentation,''
\url{https://docs.openstack.org/openstacksdk/latest/}, accessed May~2026.

\bibitem{haproxy}
HAProxy Technologies, ``HAProxy Configuration Manual,''
\url{https://www.haproxy.org/download/2.8/doc/configuration.txt},
accessed May~2026.

\bibitem{nginx_stream}
F5/NGINX, ``Module ngx\_stream\_proxy\_module,''
\url{https://nginx.org/en/docs/stream/ngx_stream_proxy_module.html},
accessed May~2026.

\bibitem{abdoc}
Apache Software Foundation, ``ab — Apache HTTP Server Benchmarking Tool,''
\url{https://httpd.apache.org/docs/2.4/programs/ab.html},
accessed May~2026.

\bibitem{ieeetempl}
IEEE, ``IEEE Conference Template,''
\url{https://www.overleaf.com/latex/templates/ieee-conference-template/grfzhhncsfqn},
accessed May~2026.

\end{thebibliography}

\end{document}
